\documentclass[9pt,t,serif,professionalfont,xcolor={table,dvipsnames},noamsthm]{beamer}
\usepackage{xr-hyper}  % for cleveref; https://tex.stackexchange.com/a/244272/121234
\PassOptionsToPackage{pdfpagemode=FullScreen}{hyperref}
\setbeamersize{text margin left = 0.5cm,
               text margin right = 0.3cm}
\mode<presentation>
{
  \usetheme{boxes}
  \setbeamercovered{invisible}
  \setbeamertemplate{navigation symbols}{} 
}
\hypersetup{colorlinks=true,linkcolor=blue} 

\usepackage{../presentation,../presentationfonts}

% load the pdfanim style, should be done after hyperref
% \usepackage{pdfanim}

% load and initialise animation
% arguments:
%   \PDFAnimLoad[options]{name}{xxx}{number}
%   - options: in this example the animation is looped
%   - name: name of the animation
%     (with this name it can be reused several times in the document)
%   - number: number of animation frames / files (n)
% the animation will be composed of files
% xxx0.pdf, xxx1.pdf ... xxx(n-1).pdf


\title{Languages}

\author{Jim Hef{}feron}
\institute{
  Mathematics and Statistics\\
  Saint Michael's College\\
  Colchester, VT\\[1ex]
  \texttt{jhefferon@smcvt.edu}
}
\date{}

% This is only inserted into the PDF information catalog.
\subject{Cardinality}

\usepackage{catchfilebetweentags} % read text from background.tex
\newcommand{\catchfilefn}{../../languages/languages.tex}
% Keep from swallowing end-of-lines
%   https://tex.stackexchange.com/a/40704/121234
\usepackage{etoolbox}
\makeatletter
\patchcmd{\CatchFBT@Fin@l}{\endlinechar\m@ne}{}
  {}{\typeout{Unsuccessful patch!}}
\makeatother

\usepackage{xr}
  \externaldocument{../../book.aux}
\usepackage{cleveref}

% \usepackage{multirow,bigstrut}
\begin{document}
\begin{frame}
  \titlepage
\end{frame}




\section{Strings}
% ------------------------------------------------------
\begin{frame}
  \frametitle{Definitions}
Recall that a \definend{symbol} (sometimes called a \definend{token}) 
is something that a machine can read and write,
and an \definend{alphabet} is a nonempty and finite set of symbols.

\ExecuteMetaData[\catchfilefn]{def:String}

We write symbols in a distinct typeface, as in \str{1} or \str{a},
because the alternative of quoting them would be 
clunky.

\ExecuteMetaData[\catchfilefn]{discussion:AlphSymbol}.
\ExecuteMetaData[\catchfilefn]{def:StringLength}
\ExecuteMetaData[\catchfilefn]{dis:OmitCommas}

\begin{example}
Natural numbers are represented by strings of digits.
The symbols are \str{0}, \str{1}, \ldots{}
and the alphabet is $\Sigma=\set{\str{0},\ldots \str{9}}$.
The string $\sigma=\str{1066}$ has length~$|\sigma|=4$. 
\end{example}

\ExecuteMetaData[\catchfilefn]{def:Bitstrings}
\end{frame}




\begin{frame}[fragile]
  \frametitle{Languages}
\ExecuteMetaData[\catchfilefn]{def:Kleenestar}

\begin{example}
Where $\Sigma=\B$, the set of bitstrings of length~$3$ is
$\B^3=\set{\str{000},\str{001}, \ldots \str{111}}$.
The set of all bitstrings is 
$\kleenestar{\B}=\set{\emptystring,\str{0},\str{1},\str{00},\ldots}$.
\end{example}

\ExecuteMetaData[\catchfilefn]{def:StringOps}

\ExecuteMetaData[\catchfilefn]{def:Lang}

\begin{example}
Let $\Sigma=\set{\str{a},\str{b}}$.
This is a language
\begin{align*}
\lang
  &=\set{\sigma\in\kleenestar{\Sigma}\suchthat\text{$\sigma=\str{a}^n\str{b}\str{a}^n$ for some $n\in\N$}}  \\
  &=\set{\str{b}, \str{aba}, \str{aabaa}, \str{aaabaaa}, \ldots}
\end{align*}
\end{example}
\end{frame}


\end{document}




